\documentclass[10pt,a4paper]{amsart}
\usepackage[margin=19mm]{geometry}
\usepackage{amsmath,amssymb,mathtools}
\usepackage[scr=rsfs,cal=euler]{mathalfa}
\usepackage{xfrac}
\usepackage[unicode,hidelinks,bookmarks=false]{hyperref}
\newcommand{\ie}{i.e.\ }
\newcommand{\cf}{cf.\ }
\newcommand{\resp}{respectively\ }
\newcommand{\Subsection}{\subsection}
\providecommand{\nobreakdash}{\nobreak\hbox{-}}
\setlength{\emergencystretch}{2em}
\multlinegap0pt
\usepackage{graphicx}
\usepackage{color}
\usepackage{cleveref}
\newcommand{\setmap}[3]{#1:#2 \mathrel{\vcenter{\mathsurround0pt
\ialign{##\crcr
		\noalign{\nointerlineskip}$\rightarrow$\crcr
		\noalign{\nointerlineskip}$\rightarrow$\crcr
		}}}%
		#3}

\newcommand{\naturals}{\mathbb{N}}
\newcommand{\real}{\mathbb{R}}
\newcommand{\realnonneg}{\mathbb{R}_{\ge 0}}
\newcommand{\realpos}{\mathbb{R}_{> 0}}
\newcommand{\until}[1]{[#1]}
\newcommand{\map}[3]{#1:#2 \rightarrow #3}
\newcommand{\qedA}{~\hfill \ensuremath{\square}}
\newcommand\scalemath[2]{\scalebox{#1}{\mbox{\ensuremath{\displaystyle #2}}}}
\newcommand{\interior}{\operatorname{int}}

\newcommand{\longthmtitle}[1]{\mbox{}{\textit{(#1):}}}
\newcommand{\setdef}[2]{\{#1 \; | \; #2\}}
\newcommand{\setdefb}[2]{\big\{#1 \; | \; #2\big\}}
\newcommand{\setdefB}[2]{\Big\{#1 \; | \; #2\Big\}}
\newcommand*{\SetSuchThat}[1][]{} % reserve the name
\newcommand*{\MvertSets}{%
    \renewcommand*\SetSuchThat[1][]{%
        \mathclose{}%
        \nonscript\;##1\vert\penalty\relpenalty\nonscript\;%
        \mathopen{}%
    }%
}
\MvertSets % default
\DeclarePairedDelimiterX \Set [2] {\lbrace}{\rbrace}
    {\,#1\SetSuchThat[\delimsize]#2\,}

\newcommand{\Ac}{\mathcal{A}}
\newcommand{\Cc}{\mathcal{C}}
\newcommand{\Dc}{\mathcal{D}}
\newcommand{\Ec}{\mathcal{E}}
\newcommand{\Fc}{\mathcal{F}}
\newcommand{\Gc}{\mathcal{G}}
\newcommand{\Hc}{\mathcal{H}}
\newcommand{\Kc}{\mathcal{K}}
\newcommand{\Pc}{\mathcal{P}}
\newcommand{\Tc}{\mathcal{T}}
\newcommand{\Uc}{\mathcal{U}}
\newcommand{\Sc}{\mathcal{S}}
\newcommand{\Vc}{\mathcal{V}}
\newcommand{\Xc}{\mathcal{X}}
\newcommand{\Yc}{\mathcal{Y}}
\newcommand{\Wc}{\mathcal{W}}
\newcommand{\Zc}{\mathcal{Z}}
\newcommand{\R}{\mathbb{R}}

\newcommand{\Ke}{\Kc^{\rm e}}

% \mydelta (paper line 96) is not reproduced: its body needs the textgreek package,
% which is absent from the runtime and is unreachable from the retained source spans.


\newcommand{\defeq}{\triangleq}

  
\newtheorem{theorem}{Theorem}
\newtheorem{lemma}{Lemma}
\newtheorem{corollary}{Corollary}
\newtheorem{proposition}{Proposition}  


\theoremstyle{definition}
\newtheorem{definition}{Definition}
\newtheorem{remarks}{Remarks}
\newtheorem{remark}{Remark}
\newtheorem{example}{Example}
\newtheorem{assumption}{Assumption}
\newtheorem{problem}{Problem}
% Bold lowercase letters
\newcommand{\ba}{\mathbf{a}}
\newcommand{\bb}{\mathbf{b}}
\newcommand{\bc}{\mathbf{c}}
\newcommand{\bd}{\mathbf{d}}
\newcommand{\be}{\mathbf{e}}
\renewcommand{\bf}{\mathbf{f}} % NOTE: use \textbf if you really want to write bold non-math text.
\newcommand{\bg}{\mathbf{g}}
\newcommand{\bh}{\mathbf{h}}
\newcommand{\bi}{\mathbf{i}}
\newcommand{\bj}{\mathbf{j}}
\newcommand{\bk}{\mathbf{k}}
\newcommand{\bl}{\mathbf{l}}
\newcommand{\bmm}{\mathbf{m}} % NOTE: \bm conflicts with the command for bold greek letters.
\newcommand{\bn}{\mathbf{n}}
\newcommand{\bo}{\mathbf{o}}
\newcommand{\bp}{\mathbf{p}}
\newcommand{\bq}{\mathbf{q}}
\newcommand{\br}{\mathbf{r}}
\newcommand{\bs}{\mathbf{s}}
\newcommand{\bt}{\mathbf{t}}
\newcommand{\bu}{\mathbf{u}}
\newcommand{\bv}{\mathbf{v}}
\newcommand{\bw}{\mathbf{w}}
\newcommand{\bx}{\mathbf{x}}
\newcommand{\by}{\mathbf{y}}
\newcommand{\bz}{\mathbf{z}}

% Bold uppercase letters
\newcommand{\bA}{\mathbf{A}}
\newcommand{\bB}{\mathbf{B}}
\newcommand{\bC}{\mathbf{C}}
\newcommand{\bD}{\mathbf{D}}
\newcommand{\bE}{\mathbf{E}}
\newcommand{\bF}{\mathbf{F}}
\newcommand{\bG}{\mathbf{G}}
\newcommand{\bH}{\mathbf{H}}
\newcommand{\bI}{\mathbf{I}}
\newcommand{\bJ}{\mathbf{J}}
\newcommand{\bK}{\mathbf{K}}
\newcommand{\bL}{\mathbf{L}}
\newcommand{\bM}{\mathbf{M}}
\newcommand{\bN}{\mathbf{N}}
\newcommand{\bO}{\mathbf{O}}
\newcommand{\bP}{\mathbf{P}}
\newcommand{\bQ}{\mathbf{Q}}
\newcommand{\bR}{\mathbf{R}}
\newcommand{\bS}{\mathbf{S}}
\newcommand{\bT}{\mathbf{T}}
\newcommand{\bU}{\mathbf{U}}
\newcommand{\bV}{\mathbf{V}}
\newcommand{\bW}{\mathbf{W}}
\newcommand{\bX}{\mathbf{X}}
\newcommand{\bY}{\mathbf{Y}}
\newcommand{\bZ}{\mathbf{Z}}

\newcommand{\bGamma}{\boldsymbol{\Gamma}}
\newcommand{\bLambda}{\boldsymbol{\Lambda}}
\newcommand{\balpha}{\boldsymbol{\alpha}}
\newcommand{\bzero}{\mathbf{0}}


\newcommand{\nom}{{\operatorname{nom}}}
\newcommand{\obs}{{\operatorname{o}}}
\newcommand{\des}{{\operatorname{d}}}
\newcommand{\sgn}{{\operatorname{sgn}}}
\newcommand{\tgt}{{\operatorname{tgt}}}
\newcommand{\fl}{{\operatorname{FL}}}
\newcommand{\predict}{{\operatorname{p}}}
\newcommand{\Lie}{\mathcal{L}}
\newcommand{\qp}{{\operatorname{QP}}}
\newcommand{\relu}{{\operatorname{ReLU}}}
\newcommand{\diag}{{\operatorname{diag}}}
\newcommand{\mineig}{{\operatorname{mineig}}}

\newcommand{\new}[1]{{\color{blue} #1}}
\newcommand{\WIP}[1]{{\color{red} #1}}
\newcommand{\Exp}{\mathbb{E}}
%\renewcommand{\new}[1]{#1}

\newcommand{\marginpo}[1]{\marginpar{\color{blue}\tiny\ttfamily#1}}

%\parskip .25ex
\renewcommand{\baselinestretch}{1}
\allowdisplaybreaks

\DeclareMathOperator*{\argmin}{argmin}
\DeclareMathOperator*{\argmax}{argmax}

%% define comment color
\newcommand{\red}{\color{red}}
\newcommand{\blue}{\color{blue}}
\newcommand{\PO}[1]{{\color{magenta} #1}}
\newcommand{\YX}[1]{{\color{blue} #1}}
\crefname{equation}{}{}
\crefname{proposition}{Proposition}{Propositions}
\title{Matrix Control Barrier Functions}
\author{Pio Ong, Yicheng Xu, Ryan M. Bena, Faryar Jabbari, and Aaron D. Ames}
\date{Source: arXiv:2508.11795v2 (CC BY 4.0); the authors' own native \LaTeX{} source, set here in the AMS \texttt{amsart} wrapper (the IEEEtran class named by the source is not present in the runtime).}
\begin{document}
\maketitle

\maketitle
\begin{abstract}
    This paper generalizes the control barrier function framework by replacing scalar-valued functions with matrix-valued ones. Specifically, we develop barrier conditions for safe sets defined by matrix inequalities---both semidefinite and indefinite. Matrix inequalities can be used to describe a richer class of safe sets, including nonsmooth ones.   The safety filters constructed from our proposed matrix control barrier functions via semidefinite programming (CBF-SDP) are shown to be continuous. Our matrix formulation naturally provides a continuous safety filter for Boolean-based control barrier functions, notably for disjunctions (OR), without relaxing the safe set. We illustrate the effectiveness of the proposed framework with applications in drone network connectivity maintenance and nonsmooth obstacle avoidance, both in simulations and hardware experiments.
\end{abstract}

\noindent\textit{Editorial scope.} The counted claim of this unit is the article's Theorem~1, \emph{Safety from exponential MCBF} (source label \texttt{thm:EMCBF-safety}), delivered with its complete author proof as the single primary proof body. Everything else retained here is uncounted support taken verbatim from the authors' own source: the abstract, Section~2 (Background on Safety Filter: notation, the barrier condition with Lemma~1, and the control barrier function background) and Section~3 up to the counted theorem, plus the \emph{Spectral Analysis} passage that follows it, including Proposition~2 with its own complete proof, which the article uses to complement the safety result. The article's Introduction (source lines 227--241), its later sections on indefinite matrix constraints, Boolean compositions, Lipschitz controllers, algorithms and experiments, and its figures are not part of the source-local core of the counted claim and are omitted; no mathematics is written, repaired or re-derived here, all mathematics is native text. Original theorem, lemma, proposition, definition, remark and equation numbers are reproduced by the authors' own \texttt{\textbackslash newtheorem} declarations (carried unchanged from their preamble), so the delivered PDF prints Definition~1, Lemma~1, Definition~2, Definition~3, Proposition~1, Remark~1, Theorem~1, Definition~4 and Proposition~2 exactly as the published article does. The source is an IEEEtran two-column document; that class is not available in the runtime, so the unit is set in the AMS \texttt{amsart} wrapper declared above, carrying the paper's own macros; the only macro not reproduced is \texttt{\textbackslash mydelta}, whose body needs the absent \texttt{textgreek} package and which is unreachable from the retained spans. Printed slips are kept literal, including the proof's cross-reference to the matrix barrier condition as ``Lemma~1'' although the label it points at is Proposition~1 (source line 392, printed the same way in the published article), the phrase ``the definition of th positive definiteness'' (source line 356), and the unbalanced parenthesis in $\partial\phi_j(\mathbf{H}(\mathbf{x}(t))$ (source line 422).
\section{Background on Safety Filter}
\subsection{Notation}
Throughout the paper, we use $\naturals$ and $\real$ for the set of natural and real numbers. For $p\in\naturals$, we denote by $\until{p}$ the set of consecutive numbers $\{1,2,\ldots,p\}$. Given a vector $\bu\in\real^m$, $\|\bu\|$ is its Euclidean norm. Given a matrix $\bM\in\real^{m\times n}$, $M_{ij}$ denotes its $(i,j)$-th entry. We use $\mathbb S^p$ for the space of symmetric real matrices of size $p\times p$. 
We use $\mathbf{1}_p\in\real^p$ for a vector of ones and $\bI_{p\times p}\in\mathbb S^p$ for the identity matrix. Given two matrices $\bA,\bB\in\real^{p\times p}$, $\bA\cdot\bB$ is their Frobenius product. Note that the Frobenius product satisfies the identity: $\bv\bv^\top\cdot \bA = \bv^\top \bA\bv$ for any $\bv\in\real^p$. A function $\map{\alpha}{\real}{\real}$ is of extended class-$\Kc$, denoted $\alpha\in\Ke$, if it is continuous, strictly increasing, and satisfies $\alpha(0)=0$. For a continuously differentiable function $\map{h}{\real^n}{\real}$, its Lie derivative along a vector field $\map{\bf}{\real^n}{\real^n}$ (or along multiple vector fields stacked in a matrix $\map{\bg}{\real^n}{\real^{n\times m}}$) is defined as $L_\bf h(\bx)=\frac{\partial h}{\partial \bx} \bf(\bx)$.



\subsection{Barrier Condition}
Consider the nonlinear autonomous system:
\begin{equation}\label{sys:auto}
    \dot \bx = \bF(\bx)
\end{equation}
where $\bx\in\real^n$ is the system state and $\map{\bF}{\real^n}{\real^n}$ is the system dynamics. For \textit{safety}, we are interested in verifying that the state trajectories $t\rightarrow\bx(t)$ remain inside a safety constraint defined as a sublevel set of the function $\map{\psi}{\real^n}{\real}$:
\begin{equation}\label{eq:safety-constraint}
    \Sc = \setdefb{\bx\in\real^n}{\psi(\bx)\geq 0},
\end{equation}
at all times. In other words, we want the set $\Sc$ to be forward invariant.
\begin{definition}
    \longthmtitle{Forward Invariance}
    \label{def:forward-invariance}
    A set $\Cc\subset\real^n$ is \textbf{forward invariant} for system~\eqref{sys:auto} if, for initial conditions $\bx_0\in\Cc$, all system trajectories $t\rightarrow\bx(t)$ remain inside the set $\Cc$ for all time $t\geq 0$. The set $\Cc$ is \textbf{safe} if it is also a subset of the safety constraint~$\Sc$.~\hfill$\diamond$
\end{definition}
The idea behind a safe set is that it provides a safe operating region, from which the system can be initialized without violating the safety constraint. In the barrier function framework, we construct the safe set using a continuously differentiable function $\map{h}{\real^n}{\real}$ as:
\begin{equation}\label{eq:safeset-scalar}
    \Cc = \setdefb{\bx\in\real^n}{h(\bx)\geq 0},
\end{equation}
typically with $h(\bx)\geq \psi(\bx),\forall \bx$ to ensure $\Cc\subset \Sc$. Forward invariance, on the other hand, can be established using the barrier condition.
\begin{lemma}\longthmtitle{Barrier Condition}
    \label{lem:BC}
    Consider the autonomous system~\eqref{sys:auto} with a continuous dynamics $\bF$ and the set $\Cc$~in \eqref{eq:safeset-scalar}. If there exists a locally Lipschitz function $\alpha\in\Ke$ such that:
    \begin{equation}\label{eq:BC-scalar}
        L_\bF h(\bx) \geq -\alpha(h(\bx))
    \end{equation}
    for all $\bx$ in an open neighborhood $\Ec\supset\Cc$, then set $\Cc$ is forward invariant for the system.~\hfill$\blacksquare$
\end{lemma}
We adopt the above version of barrier conditions from \cite{PG-JC-ME:17} because it only requires the system dynamics to be continuous, rather than locally Lipschitz.
This is made possible by requiring the barrier condition to hold over an open neighborhood set $\Ec$, instead of $\Cc$, so that the result does not need to rely on Nagumo's theorem. This distinction is particularly relevant for control systems.
\subsection{Control Barrier Function}
We consider the nonlinear control-affine system:
\begin{equation}\label{sys:ctrl-affine}
    \dot \bx = \bf(\bx)+\bg(\bx)\bu
\end{equation}
where $\bx\in\real^n$ is the system state, $\bu\in\real^m$ is the control input, $\map{\bf}{\real^n}{\real^n}$ is the system drift, and $\map{\bg}{\real^n}{\real^{n\times m}}$ is the control matrix.

For control systems, the control input $\bu$ provides additional flexibility for maintaining safety of the system. Analogous to forward invariance for autonomous system, we rely on the concept of control invariance.
\begin{definition}\longthmtitle{Control Invariance}
    A set $\Cc$ is (forward) \textbf{control invariant} for system \eqref{sys:ctrl-affine} if, for any initial condition $\bx(0)\in\Cc$, the system trajectories $t\rightarrow\bx(t)$ can be maintained inside the set $\Cc$ using some corresponding control input $t\rightarrow\bu(t)$. The set $\Cc$ is \textbf{safe} if it is also a subset of the safety constraint $\Sc$.\hfill$\diamond$
\end{definition}

Building on the concept of barrier conditions, control barrier functions are useful for verifying control invariance of a set.
\begin{definition}\longthmtitle{Control Barrier Functions}\label{def:CBF}
    A continuously differentiable function $\map{h}{\real^n}{\real}$ is called a \textbf{control barrier function} (CBF) for system~\eqref{sys:ctrl-affine} if there exists a locally Lipschitz function $\alpha\in\Ke$ such that, for each $\bx$ in the set $\Cc$ defined in~\eqref{eq:safeset-scalar}, there exists a $\bu\in\real^m$ satisfying:
    \begin{equation}\label{eq:CBC}
        \underbrace{L_\bf h(\bx)+L_\bg h(\bx)\bu}_{\dot h(\bx,\bu)} > -\alpha(h(\bx)).
    \end{equation}
\end{definition}
The idea behind CBFs is that they ensure the feasibility of designing a controller to satisfy~\eqref{eq:BC-scalar}. However, the safety result from Lemma~\ref{lem:BC} additionally requires that the controller be continuous and that the barrier condition be satisfied at all points in some neighborhood $\Ec$ outside of $\Cc$. To this end, the modern definition of a CBF uses a strict inequality, as in~\eqref{eq:CBC}, to facilitate the design of controllers that satisfy these condition.

The safety filter framework is one of the main features of CBFs. With CBFs $\{h_i\}_{i=1}^p$, we may construct an optimization-based controller to simultaneously satisfy all constraints:
\begin{align}\label{eq:CBF-QP}
    \bk(\bx)  =\argmin_{\bu\in\real^m} \quad & \|\bu-\bk_\des(\bx)\|^2                                                                  \\
             \textup{s.t.} \quad & \dot h_i(\bx,\bu) \geq -\alpha(h_i(\bx)),~\forall i\in\until{p} \nonumber
\end{align}
where $\map{\bk_\des}{\real^n}{\real^m}$ is a desired continuous controller, without considering safety. The controller~\eqref{eq:CBF-QP} effectively chooses the control input $\bu$ closest to the desired value $\bk_\des(\bx)$ while respecting the safety constraints given by the CBFs.


If the optimization is feasible\footnote{The compatibility between the multiple CBFs is an active research topic and is beyond the scope of this paper.}, the resulting controller is guaranteed continuous. This stems from the strict inequality in~\eqref{eq:CBC} that makes the optimization satisfy Slater's condition at each $\bx$, see \cite[Prop. 2.19]{RAF-PVK:96} and \cite{PM-AA-JC:25}. In addition, the strictness allows for $\bk$ to be well-defined on a neighborhood $\Ec$ of the intersection of safety constraints. As a result, safety filter works synergistically with CBFs to provide a simple construction of a continuous controller that addresses multiple safety constraints simulteneously.

There is also a nonsmooth version of barrier conditions and CBFs where the differentiability requirement on $h$ can be discarded. The details can be found in \cite{PG-JC-ME:17}. For better exposition, we will discuss them when needed later in the paper, rather than in this background section.

\begin{remark}[Lipschitz controllers]
    Local Lipschitz continuity is a desirable property for a controller. Beyond uniqueness of the closed-loop system solutions, Lipschitzness provides useful bounds and is important for system integration, e.g., cascaded or hierarchical system guarantees. The original conception on CBF~\cite{ADA-SC-ME-GN-KS-PT:19} also requires the controller to be locally Lipschitz. Despite the lack of such a guarantee, researchers have employed safety filters with great success. The state of the art results on Lipschitz guarantees are recently provided in~\cite{PM-AA-JC:25}.~\hfill$\bullet$
\end{remark}


\section{Semidefinite Matrix Constraints}
\label{sec:EMCBF}
We introduce matrix-valued barrier functions that enable a richer representation of safety constraints. Rather than using a scalar-valued function, we construct the safe set via a continuously differentiable matrix-valued function $\map{\bH}{\real^n}{\mathbb S^p}$:
\begin{equation}
    \label{eq:safeset-matrix}
    \Cc = \setdefb{\bx\in\real^n}{\bH(\bx)\succeq 0}
\end{equation}
To handle matrix barrier functions, we first introduce the following notation for the entry-wise Lie derivative matrix $\map{\bL_\bF\bH}{\real^n}{\mathbb S^p}$ along a vector field~$\bF$ as:
$$
    [\bL_\bF\bH]_{ij}(\bx) = L_\bF H_{ij}(\bx), \forall i,j\in\until{p}.
$$
Using this notation, we present a barrier condition that maintains matrix semidefiniteness as follows.
\begin{proposition}\longthmtitle{Exponential Semidefinite Matrix Barrier Condition}
    \label{prop:BC-matrix}
    Consider the autonomous system~\eqref{sys:auto} with a continuous vector field $\bF$ and the set $\Cc$ in~\eqref{eq:safeset-matrix}. If the following barrier condition holds with a positive constant $c_\alpha>0$:
    \begin{equation}
        \label{eq:BC-matrix}
        \bL_\bF\bH(\bx) \succeq -c_\alpha\bH(\bx)
    \end{equation}
    for all $\bx$ in an open neighborhood $\Ec\supset\Cc$, then set $\Cc$ is forward invariant for the system.
\end{proposition}
\begin{proof}
    For any given vector~$\bv\in\real^p$, we define an auxiliary function:
    \begin{equation}\label{eq:aux}
        \xi_\bv(\bx) = \bv^\top \bH(\bx)\bv
    \end{equation}
    Letting $t\rightarrow \bx(t)$ be any solution from an initial condition $\bx_0\in\Cc$, we have:
    \begin{align*}
        \frac{d}{dt}(\xi_\bv \circ \bx)(t) & = \bv^\top \frac{d}{dt} (\bH \circ \bx)(t) \bv \\
                                           & = \bv^\top \bL_\bF\bH(\bx(t)) \bv              \\
                                           & \geq -c_\alpha\bv^\top (\bH \circ \bx)(t) \bv  \\
                                           & = -c_\alpha (\xi_\bv \circ \bx)(t)
    \end{align*}
    where we have used the matrix inequality~\eqref{eq:BC-matrix} to derive the inequality for all time $t$ such that $\bx(t)\in\Ec$.

    Suppose there exists a solution $\bx(t)$ starting from $\bx(0)=\bx_0\in\Cc$ yet  $\bx(t^*)\not \in \Cc$ at some time $t^*\geq 0$. Then there exists $\bv^*$ such that $\xi_{\bv^*}(\bx(t^*))<0$. By continuity of the solution, there exists a time interval $[t_{\partial\Cc},t_\Ec] \subseteq [t_0,t^*]$ such that $\xi_{\bv^*}(\bx(t_{\partial\Cc}))=0$ and $\xi_{\bv^*}(\bx(t))<0$ for all $t\in(t_{\partial\Cc},t_\Ec]$ which, by the definition of th positive definiteness, indicates $\bx(t)\in\Ec\setminus \Cc, \forall t\in(t_{\partial\Cc},t_\Ec]$. However, the inequality above suggests $\xi_{\bv^*}\circ \bx$ must be increasing during time $t\in(t_{\partial\Cc},t_\Ec)$, which is a contradiction since $\xi_{\bv^*}(\bx(t_\Ec))<0=\xi_{\bv^*}(\bx(t_{\partial\Cc}))$.
\end{proof}
\cref{prop:BC-matrix} establishes a matrix  barrier condition~\eqref{eq:BC-matrix} for the semidefinite matrix safety constraint~\eqref{eq:safeset-matrix}. Analogous to the scalar case, we propose the following definition for matrix control barrier functions.
\begin{definition}\longthmtitle{Exponential MCBF}
    A continuously differentiable function $\map{\bH}{\real^n}{\mathbb S^p}$ is called an \textbf{exponential matrix control barrier function} (exponential MCBF) for system~\eqref{sys:ctrl-affine} if there exists a constant $c_\alpha> 0$ such that, for each $\bx$ in the set $\Cc$ defined in~\eqref{eq:safeset-matrix}, there exists a $\bu\in\real^m$ satisfying:
    \begin{equation}\label{eq:CBC-matrix}
        \underbrace{\bL_\bf\bH(\bx) +\sum_{i=1}^m\bL_{\bg_i}\bH(\bx)u_i}_{\dot \bH(\bx,\bu)} \succ -c_\alpha\bH(\bx).
    \end{equation}
    where $\bg_i$ is the $i$-th column of the control matrix $\bg(\bx)$.~\hfill$\diamond$
\end{definition}
The definition of MCBFs with a strictly positive definite condition helps ensure the existence of a continuous controller, establishing safety of the set $\Cc$.



\begin{theorem}\longthmtitle{Safety from exponential MCBF}

    \label{thm:EMCBF-safety}
    Consider the control-affine system~\eqref{sys:ctrl-affine} and the set  $\Cc$ in~\eqref{eq:safeset-matrix}. If $\bH$ is an exponential MCBF, then set $\Cc$ is control invariant. In particular, the CBF-based semidefinite programming (CBF-SDP) safety filter:
    \begin{align}\label{eq:CBF-SDP}
        \bk(\bx) = \argmin_{\bu\in\real^m} \quad &\|\bu-\bk_\des(\bx)\|^2\hfill\nonumber      \\
                 \textup{s.t.} \quad & \dot \bH(\bx,\bu) \succeq -c_\alpha\bH(\bx)
    \end{align}
    is continuous for all $\bx$ in some neighborhood $\Ec$ of $\Cc$. Consequently, the state-feedback $\bu=\bk(\bx)$ renders set $\Cc$ forward invariant for the closed-loop system.
\end{theorem}
\begin{proof}

    Define the set-valued map:
    \begin{align*}
        \Uc(\bx) & = \setdefb{\bu\in\real^m}{\dot\bH(\bx,\bu) +c_\alpha\bH(\bx) \succeq \bzero}        \\
                 & =\setdefb{\bu\in\real^m}{\phi_1\big(\dot\bH(\bx,\bu) +c_\alpha\bH(\bx)\big)\geq 0}.
    \end{align*}
    where $\phi_1$ is the function returning the minimum eigenvalue of a given matrix. Since the matrices are entry-wise continuously differentiable and the eigenvalue function~$\phi_1$ is globally Lipschitz (see e.g., \cite[Thm. 1]{MDS-JC:09} and \cite[Thm. 2.4]{ASL:96}), the resulting function $\phi_1\circ(\dot \bH+c_\alpha\bH)$ is continuous. As such, not only does the map have a nonempty interior for $\bx\in\Cc$ from the definition of MCBF, but it also has a nonempty interior on some neighborhood $\Ec\supset \Cc$ from continuity. Hence, we conclude the map $\Uc$ is lower semicontinuous on $\Ec$ because it has a nonempty interior, is convex-valued, and the function $\phi_1\circ(\dot \bH+c_\alpha\bH)$ defining it is continuous, see \cite[Lem. 5.2]{GS:18}. After substituting $\Delta \bu = \bu-\bk_\des(\bx)$ as a new variable, we have a minimal selection formulation:
        \begin{align*}
            \Delta\bk(\bx) = \argmin_{\bu\in\real^m} \quad &\|\Delta\bu\|^2\hfill\nonumber                                          \\
                          \textup{s.t.} \quad & \phi_1(\dot \bH\big(\bx,\Delta\bu+\bk_\des(\bx))+c_\alpha\bH(\bx)\big) \geq 0.
        \end{align*}
    Here, the constraint remains lower semicontinuous, convex-valued, and closed-valued.  Therefore, from \cite[Prop. 2.19]{RAF-PVK:96}, its minimal selection on a Euclidean space is continuous for all $\bx\in\Ec$, implying $\bk(\bx)=\Delta\bk(\bx)+\bk_\des(\bx)$ is also continuous, as desired.

    Forward invariance of the closed-loop system $\bF(\bx) = \bf(\bx)+\bg(\bx)\bk(\bx)$ follows from Lemma~\ref{prop:BC-matrix}, and control invariance follows from the existence of the control signal $t\rightarrow\bu(t)=\bk(\bx(t))$, concluding the proof.
\end{proof}
Theorem~\ref{thm:EMCBF-safety} formally states the safety result established by an exponential MCBF, proposing specifically the CBF-SDP safety filter~\eqref{eq:CBF-SDP} as one of the possible controllers. Beyond safety guarantees, CBFs typically offer useful bounds on the evolution of safety. To complement our results, we next provide  the bounds associated with MCBF-based controllers.



\subsection{Spectral Analysis}
For the standard CBF condition~\eqref{eq:BC-scalar}, we can derive an exponential bound on the evolution of $h$ along the trajectory. This occurs when $\alpha(h(\bx))= c_\alpha h(\bx)$ is chosen as a linear function with $c_\alpha >0$. That is, we have $h(\bx(t))\geq h(\bx_0)e^{-c_\alpha t}$ from an initial condition $\bx_0\in\real^n$. Incidentally, we also refer to our MCBF in~\eqref{eq:CBC-matrix} as exponential, despite $\bH$ not being scalar-valued. This is motivated by how the evolution of its eigenvalues can be similarly exponentially bounded.

For each $\bx\in\real^n$, $\bH(\bx)$ is a symmetric matrix and thus has real eigenvalues. Formally, we use $\map{\phi_j}{\mathbb S^p}{\real}$ to denote the function mapping a symmetric matrix $\bH$ to its $j$-th smallest eigenvalue. At the same time, we use $\map{\lambda_j}{\real^n}{\real}$ for the state-dependent eigenvalue function through a function composition: $\lambda_j(\bx)=\phi_j(\bH(\bx))$. Without loss of generality, we define $\{\lambda_j\}_{j=1}^p$ so that they are in an ascending order:
\begin{equation}
    \lambda_1(\bx)\leq \dots \leq \lambda_p(\bx),~\forall \bx\in\real^n.
\end{equation}
We note importantly that the function $\{\lambda_j\}_{j=1}^p$ is not differentiable at points where the eigenvalue  is not simple. Therefore, we rely on nonsmooth analysis to bound its rate of change along an autonomous system~\eqref{sys:auto} as:
$$
    \frac{d}{dt}(\lambda_j\circ \bx)(t) \in \partial \phi_j (\bH(\bx(t))) \cdot \bL_\bF\bH(\bx(t)),  ~a.e.~t\geq0,
$$
with its weak set-valued Lie derivative~\cite[Lem. 1 and Rmk. 1]{PG-JC-ME:17}, constructed by the generalized gradient set $\partial\phi_j$. Note that the left hand side is also guaranteed to be absolutely continuous. For the right hand side, we have used the nonsmooth chain rule~\cite[Thm. 2.3.10]{FHC:83} for the weak set-valued Lie derivative since we know the generalized gradient of an eigenvalue of a symmetric matrix function, cf.~\cite{MDS-JC:09}:
$$
    \partial\phi_j(\bH) = \text{co}\setdefb{\bv\bv^\top \in \mathbb{S}^p}{\bH\bv=\phi_j(\bH)\bv,~\|\bv\| = 1}.
$$
Since the vectors $\bv$'s defining in the set above are essentially the normalized eigenvectors associated with $\phi_j$, we have the following result. 

\begin{proposition}\longthmtitle{Exponential Bound on Eigenvalues}\label{prop:exp_bound}
    Consider the autonomous system~\eqref{sys:auto} with continuous dynamics $\bF$ and the set $\Cc$ in~\eqref{eq:safeset-matrix}. If the matrix barrier condition~\eqref{eq:BC-matrix} holds, then the eigenvalues can be bounded as:
    \begin{equation}\label{eq:exp_bound}
        \lambda_i(\bx(t)) \geq \lambda_i(\bx_0)e^{-c_\alpha t},~\forall i\in\until p,
    \end{equation}
    at almost every time $t\geq0$,  along any Carath\'eodory solution starting from $\bx_0\in\Cc$.
\end{proposition}
\begin{proof}
    Let $\bA(t)\in\partial\phi_j(\bH(\bx(t))$ be such that
    $$
        \frac{d}{dt}(\lambda_j\circ \bx)(t) = \bA(t)\cdot \dot \bH(\bx(t),\bk(\bx(t)))
    $$
    at almost every time $t\geq 0$. Since $\bA$ is an element in the convex hull, there may not be an eigenvector $\bv$ such that $\bA=\bv\bv^\top$. Therefore, we use the Carath\'eodory theorem of convex hulls~\cite[Thm. 17.1]{RTR:70} to deduce the existence of  vectors $\{\bv_s(t)\}_{s=1}^{(p^2+p)/2}$ such that
    $$
        \bA(t) = \sum_{s=1}^{(p^2+p)/2} a_s(t) \bv_s(t)\bv_s(t)^\top
    $$
    where $\sum a_s(t) = 1$ and each $\bv_s$ is a normalized eigenvector associated with $\phi_j(\bH(\bx))$. Using this fact, we bound:
    \begin{align*}
        \bA\cdot \bL_\bF\bH(\bx) & = \Big(\sum_{s=1}^{(p^2+p)/2} a_s \bv_s\bv_s^\top\Big) \cdot \bL_\bF\bH(\bx) \\
                                 & = \sum_{s=1}^{(p^2+p)/2} a_s \big(\bv_s\bv_s^\top \cdot \bL_\bF\bH(\bx)\big) \\
                                 & = \sum_{s=1}^{(p^2+p)/2} a_s \big(\bv_s^\top \bL_\bF\bH(\bx)\bv_s\big)       \\
                                 & \geq  -c_\alpha\sum_{s=1}^{(p^2+p)/2} a_s \big(\bv_s^\top \bH(\bx)\bv_s\big) \\
                                 & = -c_\alpha\sum_{s=1}^{(p^2+p)/2} a_s \lambda_j(\bx)                   = -c_\alpha \lambda_j(\bx)
    \end{align*}
    where we have dropped the dependencies on $t$ for compactness of the presentation. Hence, the comparison lemma ensures the evolution bound~\eqref{eq:exp_bound}, concluding the proof.
\end{proof}
Proposition~\ref{prop:exp_bound} shows that the proposed matrix barrier condition~\eqref{eq:BC-matrix} enforces bounds on all eigenvalues of $\bH$. The bounds suggest a degree of conservatism in the approach, since ensuring only the smallest eigenvalue $\lambda_1$ nonnegative is sufficient (and necessary) to keep $\bH$ positive semidefinite. Nevertheless, it is typically difficult to handle the nonsmooth nature associated with eigenvalues, and our approach offers a simple way to address it. We further expand on this point by focusing our discussion on diagonal matrices.


\begin{thebibliography}{99}
\bibitem[1]{ADA-SC-ME-GN-KS-PT:19} A.~D. Ames, S.~Coogan, M.~Egerstedt, G.~Notomista, K.~Sreenath, and P.~Tabuada, ``Control barrier functions: Theory and applications,'' in {\em {E}uropean {C}ontrol {C}onference}, (Naples, Italy), pp.~3420--3431, June 2019.
\bibitem[9]{RAF-PVK:96} R.~A. Freeman and P.~V. Kototovic, {\em Robust Nonlinear Control Design: State-space and {L}yapunov Techniques}. \newblock Boston, MA: Birkh{\"a}user, 1996.
\bibitem[10]{PM-AA-JC:25} P.~Mestres, A.~Allibhoy, and J.~Cortés, ``Regularity properties of optimization-based controllers,'' {\em European Journal of Control}, vol.~81, p.~101098, 2025.
\bibitem[12]{PG-JC-ME:17} P.~Glotfelter, J.~Cort\'es, and M.~Egerstedt, ``Nonsmooth barrier functions with applications to multi-robot systems,'' {\em IEEE Control Systems Letters}, vol.~1, no.~2, pp.~310--315, 2017.
\bibitem[16]{MDS-JC:09} M.~D. Schuresko and J.~Cort{\'e}s, ``Distributed motion constraints for algebraic connectivity of robotic networks,'' {\em Journal of Intelligent and Robotic Systems}, vol.~56, no.~1, pp.~99--126, 2009.
\bibitem[21]{ASL:96} A.~S. Lewis, ``Group invariance and convex matrix analysis,'' {\em SIAM Journal on Matrix Analysis and Applications}, vol.~17, no.~4, pp.~927--949, 1996.
\bibitem[22]{GS:18} G.~Still, ``Lectures on parametric optimization: An introduction,'' {\em Optimization Online}, 2018.
\bibitem[23]{FHC:83} F.~H. Clarke, {\em Optimization and Nonsmooth Analysis}. \newblock Canadian Mathematical Society Series of Monographs and Advanced Texts, New York: Wiley, 1983.
\bibitem[24]{RTR:70} R.~T. Rockafellar, {\em Convex Analysis}. \newblock Princeton, NJ: Princeton University Press, 1970.
\end{thebibliography}
\end{document}
